\subsection{A $p$-separation Markov property}

We will ignore nullset subtleties for now, and use the easy quantifiers (for all values of exogenous variables).
We will also omit input nodes for now.

We define a graphical operation that we call ``teleportation'':
\begin{Def}\label{def:teleportation_graph}
  Let $G$ be a DMG with nodes $V$.
  Write $I(G) := \{ (i,j) \in V^2 : i \tuh j \in G, i \in \Sc_G(j) \}$ be the set of node pairs that lie within a non-trivial strongly connected component of $G$ and are connected by a directed edge.
  We define the \emph{teleportation of $G$} as the DMG $G_{tp}$ with
  \begin{itemize}
    \item nodes $V \dcup \{s_{(i,j)} : (i,j) \in I(G)\}$;
    \item directed edges 
      $$\{ i \tuh j : i \tuh j \in G, i \notin \Sc_G(j) \} \cup \{ i \tuh s_{(i,j)} : (i,j) \in I(G) \}$$
    \item bidirected edges
      $$\{ i \huh j : i \huh j \in G \} \cup \{ s_{(i,j)} \huh j : (i,j) \in I(G) \}.$$
  \end{itemize}
\end{Def}
\begin{Prp}
  Let $G$ be a DMG with nodes $V$.
  Then $G_{tp}$ is acyclic.
\end{Prp}
\begin{proof}
  \Joris{TODO}
\end{proof}
We define an analogous operation for SCMs:
\begin{Def}\label{def:teleportation_scm}
  Let $M = \lt V,W,\CX,P,f \rt$ be an SCM.
  We call an SCM $\tilde{M} = \lt \tilde{V},\tilde{W},\tilde{\CX},\tilde{P},\tilde{f} \rt$ a \emph{teleportation of $M$} if:
  \begin{itemize}
    \item $\tilde{V} := V \dcup S$, where $S$ is as in Definition~\ref{def:teleportation_graph} (applied to $G(M)$);
    \item $\tilde{W} := W \dcup U$, where $U := \{u_{(i,j)} : (i,j) \in I(G)\}$;
    \item $\CX := \CX \times \prod_{s \in S} \CX_s \times \prod_{u \in U} \CX_u$,
      where $\CX_s := \{0,1\}$ for each $s \in S$, and $\CX_{u_{(i,j)}} := \CX_i$ for each $u_{(i,j)} \in U$;
    \item $\tilde{P}(X_W) = P(X_W)$;
    \item \Joris{getting the notation right is a little challenging, we should probably draw inspiration from the node-splitting operation}
      The causal mechanism $\tilde{f}$ is obtained from the causal mechanism $f$ in the following way.
      For the new endogenous variables $S$, we require:
        $$\tilde f_{s_{(i,j)}}(x_V,x_S,x_W,x_U) = \ind[x_i = x_{u_{ij}}].$$
        For the original endogenous variables in $V$, let $I_j := \{i : (i,j) \in I(G)\}$ (the former parents of $j$ in $\Sc_G(j)$), and $U_j := \{u_{(i,j)} \in I(G)\}$ (the associated newly introduced exogenous random variables).
      Then we require:
        $$\tilde f_j(x_V,x_S,x_W,x_U) = f_j((x_{V \sm I_j},x_{U_j}),x_W)$$
  \end{itemize}
  So, in words, we introduce a binary endogenous variable $s_{(i,j)}$ for each edge in $I(G)$, and an exogenous random variable $u_{(i,j)}$ which takes values in the same space as the endogenous variable $i$.
  Also, we keep the causal mechanisms of $M$ but ``rewire'' them by replacing each parent $i$ of $j$ by the corresponding $u_{(i,j)}$.
  We shall write $M_{tp}$ for an (arbitrary) teleportation of $M$.
  Note that the conditional distribution $P_{M_{tp}}(X_U \given X_W)$ is not constrained in any way.
\end{Def}
The graphical and SCM teleportation operations are compatible.
\begin{Prp}
  For an SCM $M$ with graph $G(M)$, we have that
  $$G(M_{tp}) = (G(M))_{tp}$$
  for any teleportation $M_{tp}$ of $M$.
  \Joris{Here a problem occurs: we should not weaken the quantifier for the dummy exogenous variables $U$ to ``for almost all''. 
  Hence, measure theoretically it seems better to treat them like exogenous input variables. 
  This also suggests that we don't need to put a distribution on those.
  However, it then conflicts with the assumption that endogenous variables cannot cause exogenous input variables.
  So the conclusion is that our current formal framework is simply not flexible/general enough to properly accommodate this object.
  For now, we will pragmatically ignore this problem.}
\end{Prp}
We can now define the notion of $p$-separation with the help of the graphical teleportation.
\begin{Def}\label{def:p-separation}
  Let $G$ be a DMG with nodes $V$.
  For $A,B,C \ins V$, we define that \emph{$A$ is $p$-separated from $B$ given $C$} as:
  $$A \Perp_{G}^p B \given C := A \Perp_{G_{tp}}^d B \given C \dcup S$$
  where $S$ are the selection nodes introduced in $G_{tp}$.
\end{Def}
\begin{Prp}
  $p$-separation satisfies the separoid axioms (and the other usual axioms for $d$-separation).
\end{Prp}
\begin{proof}
  Easy to see: $p$-separation only replaces $G$ with $G_{tp}$ and $C$ with $C \dcup S$, and in the axioms the conditioning set can only grow, but it never shrinks (and hence the conditioning set of the implied separation will still contain $S$).
\end{proof}

\begin{Rem}
  A simple way to derive a $p$-separation Markov property in the special case when all variables are discrete is the following.
  
  For $A,B,C \ins V \cup W$,
  $$A \Perp_{G^+(M)}^p B \given C \implies A \Perp_{(G^+(M))_{tp}}^d B \given C \dcup S \implies A \Perp_{G^+(M_{tp})}^d B \given C \dcup S.$$
  We can apply the d-separation Markov property to $M_{tp}$ (because it is acyclic).
  This states that for all subsets $A,B,C \ins \tilde{V} \cup \tilde{W} = V \cup S \cup W \cup U_{tp}$,
  $$A \Perp_{G^+(M_{tp})}^d B \given C \implies X_A \Indep_{P_{M_{tp}}} X_B \given X_C.$$
  Hence we conclude: for $A,B,C \ins V \cup W$,
  $$A \Perp_{G^+(M)}^p B \given C \implies X_A \Indep_{P_{M_{tp}}} X_B \given X_C,X_S.$$

  If $M$ is uniquely solvable, and if $P_{M_{tp}}(X_U)$ has full support, then $P_{M_{tp}}(X_S=1) > 0$
  and $P_{M_{tp}}(X_V,X_W \given X_S=1) = P_M(X_V,X_W)$.
  Hence, under those assumptions:
  \begin{align*}
    X_A \Indep_{P_{M_{tp}}(X_V,X_W,X_S)} X_B \given X_C,X_S 
    & \implies X_A \Indep_{P_{M_{tp}}(X_V,X_W,X_S)} X_B \given X_C,X_S=1 \\
    & \implies X_A \Indep_{P_{M_{tp}}(X_V,X_W|X_S=1)} X_B \given X_C \\
    & \implies X_A \Indep_{P_M(X_V,X_W)} X_B \given X_C
  \end{align*}
\end{Rem}

For continuous variables, however, this strategy runs into the problem that then $P_{M_{tp}}(X_S=1)$ may vanish, so we cannot condition on that event.
The remainder sketches a way out.\footnote{After having worked this out, I noticed that for convenience I worked with mathematical structures that are
very similar to SCMs, and the terminology I chose reflects this similarity. However, these structures come with a very non-standard causal interpretation,
so in order to avoid confusion it might have been better to pick other names.}

\begin{Def}\label{def:sscm}
  Given a tuple $M = \lt V \dcup S,W \dcup U,\CX,P,\tilde f \rt$ that is ALMOST an SCM except that we only specify the marginal $P_M(X_W) = \bigotimes_{w \in W} P_M(X_w)$ instead of the joint $P_M(X_W,X_U)$,
  and such that $\Ch^{G^+(M)}(S) = \emptyset$, and a value $s \in \CX_S$, we call the tuple $(M,s)$ an sSCM.

  A distribution $\Prb(X_V,X_S,X_W,X_U)$ is called a \emph{solution of the sSCM $(M,s)$} if
  \begin{enumerate}
    \item $\Prb(X_S=s)=1$, and
    \item $\Prb(X_{V \cup S} = \tilde f(X_V,X_S,X_W,X_U))=1$, and
    \item $\Prb(X_W) = P_M(X_W)$.
  \end{enumerate}

  The \emph{graph} of sSCM $(M,s)$ is defined as the graph $G^+(M)$ of $M$ (after complementing with an arbitrary independent factorizing exogenous distribution on $\CX_U$).
\end{Def}

\begin{Prp}
  If $M$ is uniquely solvable, then $(M_{tp},1)$ has a solution.
\end{Prp}
\begin{proof}
  \Joris{TODO. Essentially, consider the (unique) solution function $g : \CX_W \to \CX_V$ of $M$.
  We can use it to construct a ``solution function'' $\tilde{g} : \CX_W \to \CX_{V \cup S \cup U} : x_w \mapsto (g(x_W),1,(g_i(x_W) : (i,j) \in I(G)))$ for $(M_{tp},1)$.
  Now take the push-forward $\tilde g_\star P(X_W)$.}
\end{proof}

\begin{Rem}
  For a node $v \in V$ such that $\Ch^{G^+(M)}(v) = \emptyset$, we can marginalize out $v$ to obtain an sSCM
  $$M_{\sm v} = \lt (V \sm \{v\}) \dcup S, W \dcup U, \CX_{(V \sm \{v\}) \dcup S \dcup W \dcup U}, P, \tilde f_{(V \sm \{v\}) \dcup S} \rt.$$
  If $(M,s)$ has solution $\Prb$, then $(M_{\sm v},s)$ has solution $\Prb_v := \Prb(X_{V\sm v},X_S,X_W,X_U)$ (the marginal distribution of $\Prb$).
  If $(M,s)$ has graph $G^+(M)$, then $(M_{\sm v},s)$ has graph $(G^+(M))^{\sm v}$ (the marginalized graph of $M$).
\end{Rem}

We use the separoid axioms, Theorem~\ref{thm:sym-sep-ax-cond-ind}.
\begin{Thm}[Conditional Global Markov property for Acyclic sSCMs]\label{thm:gmp_acyclic_selection_scms}
  Given an sSCM $(M,s)$ as in Definition~\ref{def:sscm}, that is acyclic with solution $P$ and graph $G = G^+(M)$.
  Then for all subsets $A,B,C \ins V \cup W$,
  $$A \Perp_G^d B \given C \dcup S \implies X_A \Indep_{P} X_B \given X_C.$$
\end{Thm}
\begin{proof}
  Induction start: $V=\emptyset$, and $W,S,U$ arbitrary.
  This means that $A,B,C \ins W$.
  By Lemma~\ref{lem:d-separation-markov-property}(i), we have 
    $$A \cap B \ins C \cup S.$$
  Since $S \cap W = \emptyset$, we conclude that
    $$A \cap B \ins C.$$
  Because $P$ is a solution of sSCM $(M,s)$, $P(X_W) = \bigotimes_{w \in W} P(X_w)$ and hence
    $$X_A \Indep_P X_B \given X_C.$$

  We now assume that the claim holds for $\#V < n$ and arbitrary $W,S,U$.
  We consider the case $\#V=n>0$.

  By Lemma~\ref{lem:d-separation-markov-property}(ii), $A\sm (C \cup S)$, $B \sm (C \cup S)$ and $C \cup S$ are pairwise disjoint.
  Because $S \cap W = \emptyset$, we conclude from this that
    $$A' := A \sm (C \cup \{v\}), \quad B' := B \sm (C \cup \{v\}), \quad C' := C \sm \{v\}, \quad \{v\}, \quad S$$
  are pairwise disjoint.

  Let $v$ be a node in $V$ such that $\Ch^G(v) = \emptyset$, which exists because we assumed $M$ (and hence $G$) to be ayclic.
  Write $D := \Pa^G(v)$.
  Note that $D \ins V \cup W \cup U$, but $D \cap S = \emptyset$ because $\Ch^G(S) = \emptyset$ by assumption.
  Since by assumption
  $$P(X_v = f_v(X_V,X_S,X_W,X_U)) = 1,$$
  we get the independence (local Markov property)
  \begin{equation}\label{eqn:blu}
    X_v \Indep_P X_{(V \sm \{v\}) \cup W} \given X_D.\tag{\#4}
  \end{equation}

  In the following we will distinguish between the 4 possible cases: 
  $v \in A \sm C$, $v \in B \sm C$, $v \in C$, $v \in V \sm (A \cup B \cup C)$.

  \begin{enumerate}[label=(\roman*)]
    \item $v \in V \sm (A \cup B \cup C)$. The claim follows from the induction hypothesis applied to $(M_{\sm v},s)$:
    \begin{align*}
      A \Perp^d_G B \given C \dcup S
      & \xRightarrow{\text{decomposition}}  & A' \Perp^d_G B' \given C \dcup S \\
      & \xRightarrow{\text{marginalization}} &  A' \Perp^d_{G^{\sm v}}  B' \given C \dcup S\\
      & \xRightarrow{\text{induction}} & X_{A'} \Indep_{P_{\sm v}} X_{B'} \given X_C\\
      & \xRightarrow{\text{marginalization}} & X_{A'} \Indep_{P} X_{B'} \given X_C.
    \end{align*}

    \item $v \in A\sm C$.
    We then have the implications:
    \begin{align*}
      A \Perp^d_G  B \given C \dcup S
      & \xRightarrow{\text{decomposition}}  & A' \Perp^d_G B' \given C \dcup S \\
      & \xRightarrow{\text{marginalization}} &  A' \Perp^d_{G^{\sm v}}  B' \given C \dcup S\\
      & \xRightarrow{\text{induction}} & X_{A'} \Indep_{P_{\sm v}}  X_{B'} \given X_C \\
      & \xRightarrow{\text{marginalization}} & X_{A'} \Indep_{P} X_{B'} \given X_C. \tag{\#1}
    \end{align*}
    On the other hand we have:
    \begin{align*}
      A \Perp^d_G  B \given C \dcup S 
      & \xRightarrow{\text{decomposition}}  & A \Perp^d_G   B' \given C \dcup S   \\
      & \xRightarrow{\text{weak union}}  & \{v\} \Perp^d_G  B' \given A' \dcup C \dcup S    \\
      & \xRightarrow{\text{Lemma~\ref{lem:d-separation-markov-property}(iii)}}  & D  \Perp^d_G   B' \given A' \dcup C \dcup S \\ 
      & \xRightarrow{\text{induction}}  &   X_D  \Indep_{P}  X_{B'} \given X_{A'},X_C \tag{\#2}
    \end{align*}
    \Joris{Oops: mistake. The induction step is invalid if $D \cap U \ne \emptyset$.}
    The local Markov property \eqref{eqn:blu} for $v$ now implies:
    \begin{align*}
    &&X_v \Indep_{P}  X_{(V \cup W) \sm \{v\}} \given X_D\\ 
    &\xRightarrow{\text{decomposition}} & 
    X_v \Indep_{P}  X_{A'}, X_{B'}, X_C \given X_D \\
    & \xRightarrow{\text{weak union}}  & 
    X_v \Indep_{P} X_{B'} \given X_{A'}, X_C, X_D \\ 
    & \xRightarrow{\text{contraction}\eqref{eqn:bla}}  & X_v, X_D  \Indep_{P}  X_{B'} \given X_{A'}, X_C \\
    & \xRightarrow{\text{decomposition}}  & X_v \Indep_{P}  X_{B'} \given X_{A'}, X_C\\
    & \xRightarrow{\text{contraction} \eqref{eqn:bli}}  &  X_{A'},X_v \Indep_{P}  X_{B'} \given X_C\\
    & \xRightarrow{\text{redundancy}}  & X_{A'}, X_v \Indep_{P}  X_B \given X_C.
     \tag{\#3}
    \end{align*}
    \Joris{The last step is not entirely clear. $B' = B \sm C$. So $B \ins B' \cup C$.
    It appears that we need to combine redundancy and contraction, but then I am missing something which appears true but not part of the separoid axioms:
      $$X \Indep Y \given C \implies X \Indep Y \given C,\varphi(C)$$
            }
    From redundancy we get the implications:
    \begin{align*}
      &&X_{A\sm C} \Indep_{P} X_B \given X_{A'},X_v,X_C\\
    & \xRightarrow{\text{contraction} \eqref{eqn:bloop}}  &  
      X_{A'},X_v,X_{A\sm C} \Indep_{P}  X_B \given X_C\\
    & \xRightarrow{\text{decomposition}}  &    
    X_A \Indep_{P}  X_B \given X_C. 
    \end{align*}
    This shows the claim in the first case.

  \item $v \in B\sm C$. This follows analogously to the previous case (by exchanging the roles of $A$ and $B$).

  \item $v \in C$.
    We get the implications:
    \begin{align*}
    A  \Perp^d_G B \given C \dcup S
      & \xRightarrow{\text{decomposition}}  & A' \Perp^d_G  B' \given C'\dcup\{v\} \dcup S \\
      & \xRightarrow{\text{Lemma~\ref{lem:d-separation-markov-property}(iv)}} &  A' \dcup \{v\} \Perp^d_G  B' \given C' \dcup S \qquad \lor \qquad A' \Perp^d_G B' \dcup \{v\} \given C' \dcup S \\
      & \xRightarrow{\text{cases (i),(ii)}} &  X_{A'},X_v \Indep_{P}  X_{B'} \given X_{C'} \quad \lor \quad X_{A'} \Indep_{P}  X_{B'},X_v \given X_{C'} \\
      & \xRightarrow{\text{weak union}} &  X_{A'} \Indep_{P}  X_{B'} \given X_{C} \quad \lor \quad X_{A'} \Indep_{P}  X_{B'} \given X_{C} \\
      & \xRightarrow{\text{redundancy}} &  X_A \Indep_{P}  X_B \given X_{C} \quad \lor \quad X_A \Indep_{P}  X_B \given X_{C}
    \end{align*}
    This shows the claim in case (iii).
\end{enumerate}
\end{proof}
We made use of the following consequences of $d$-separation.
\begin{Lem}\label{lem:d-separation-markov-property}
  Let $G$ be a DMG with nodes $V$, and let $A,B,C \ins V$ such that $A \Perp_G^d B \given C$.
  Then:
    $$A \Perp_G^d B \given C$$
  implies:
  \begin{enumerate}[label=\roman*)]
    \item $A \cap B \ins C$.
    \item $A\sm C$, $B \sm C$ and $C$ are pairwise disjoint.
    \item If $v \in A \sm C$ is such that $\Ch_G(v) = \emptyset$, then
      $$\{v\} \Perp^d_G B' \given A' \dcup C \implies \Pa^G(v) \Perp^d_G B' \given A' \dcup C.$$
    \item If $v \in C$ is such that $\Ch_G(v) = \emptyset$, one of the following statements holds:
      \[  A' \dcup \{v\} \Perp^d_G  B' \given C' \qquad \lor \qquad A' \Perp^d_G B' \dcup \{v\} \given C'. \]
      where writing $A' := A \sm C$, $B' := B \sm C$, $C' := C \sm \{v\}$. 
  \end{enumerate}
\end{Lem}
\begin{proof}
  \begin{enumerate}[label=\roman*)]
    \item 
      Otherwise, a trivial path between $A$ and $B$ would be open.
    \item 
      This follows immediately from the first point.
    \item
      This holds since every $(A' \dcup C)$-open walk $ w \sus \cdots$ from a  $w \in \Pa^G(v)$ to $B'$ 
      extends to an $(A' \dcup C)$-open walk from $v$ to $B'$ via  $v \hut w \sus \cdots$, as $w$ stays a 
      non-collider in the extended walk (not in $A' \dcup C$) and $v \notin A' \dcup C$.
    \item
      By decomposition, and using the previous point,
      $$A \Perp_G^d B \given C$$
      implies
      $$A' \Perp_G^d B' \given C' \dcup \{v\}.$$
      Assume the contrary of the claim:
      \[  A' \dcup \{v\} \nPerp^d_G  B' \given C' \qquad \land \qquad A' \nPerp^d_G B' \dcup \{v\} \given C'. \]
      So there exist shortest $C'$-open walks $\pi_1$ and $\pi_2$ in $G$ such that all colliders are in $C'$:
      \[ \pi_1:\quad  A' \cup \{v\}\ni u_0 \sus \cdots \sus u_k \in B',  \]
      and:
      \[ \pi_2:\quad  A' \ni w_0 \sus \cdots \sus w_m \in (B' \dcup \{v\}).  \]
      So all non-colliders of $\pi_1$ and $\pi_2$ are outside of $C'$.
      Since we consider shortest walks and $v \notin C'$ at most an end node of $\pi_1$ and $\pi_2$ could be equal to $v$.
      Otherwise one could shorten the walk.
      Then note that $v \notin A'$ and $v \notin B'$, thus: $u_k \neq v$ and $w_0 \neq v$.
      If now $\pi_i$ does not contain $v$ as an (end) node, then $\pi_i$ would be $(C'\dcup \{v\})$-open,
      which is a contradiction to the assumption:
      \[A' \Perp^d_G  B' \given C'\dcup\{v\}.\]
      So we can assume that the other end nodes equal $v$, i.e.: $u_0=v$ and $w_m =v$.
      Furthermore, both $\pi_1$ and $\pi_2$ are non-trivial walks, since $u_0 \neq u_k$ and $w_0 \neq w_m$.
      %Otherwise $v=u_k \in B'$ or $v=w_0 \in A'$, both of which are not true.
      %So $k,m \ge 1$.\\
      Since $v$ is childless and $k,m \ge 1$ we have that the $\pi_i$ are of the forms:
      \[ \pi_1:\quad   v \hut u_1 \sus \cdots \sus u_k,  \]
      and:
      \[ \pi_2:\quad  w_0 \sus \cdots \sus w_{m-1} \tuh v,  \]
      with $u_1, w_{m-1} \in D=\Pa^G(v)$. Then the following walk:
      \[ A' \ni  w_0 \sus \cdots \sus w_{m-1} \tuh v \hut u_1 \sus \cdots \sus u_k \in B',  \]
      is a $(C' \dcup \{v\})$-open walk from $A'$ to $B'$, in contradiction to:
      \[A' \Perp^d_G B' \given C'\dcup\{v\}.\]
\end{enumerate}
\end{proof}

Now we can prove a $p$-separation Markov property for uniquely solvable SCMs.
\begin{Cor}
  Let $M$ be a uniquely solvable SCM.
  For $A,B,C \ins V \cup W$,
  $$A \Perp_{G^+(M)}^p B \given C \implies X_A \Indep_{P_M(X_V,X_W)} X_B \given X_C.$$
\end{Cor}
\begin{proof}
  Pick a teleportation $M_{tp}$.
  The sSCM $(M_{tp},1)$ is acyclic.
  Because $M$ is uniquely solvable, the sSCM $(M_{tp},1)$ has a solution $P_{M_{tp}}$.
  For $A,B,C \ins V \cup W$,
  $$A \Perp_{G^+(M)}^p B \given C \implies A \Perp_{(G^+(M))_{tp}}^d B \given C \dcup S \implies A \Perp_{G^+(M_{tp})}^d B \given C \dcup S.$$
  We can now apply the conditional global Markov property Theorem~\ref{thm:gmp_acyclic_selection_scms} to the acyclic sSCM $(M_{tp},1)$ with solution $P_{M_{tp}}$, which gives the implication:
  $$A \Perp_{G^+(M_{tp})}^d B \given C \dcup S \implies X_A \Indep_{P_{M_{tp}}} X_B \given X_C.$$
  Because
  $$P_{M_{tp}}(X_V,X_W) = P_M(X_V,X_W),$$
  the claim follows.
\end{proof}
